% Section 2, part 3 of 8: the corollary of prop:set-cyclicLR; points; Frobenius data.

\begin{cor}[Eigenvalues, idempotents and the fifth-order equation]
\label{cor:set-cyclicPF5}
Assume (M) and (G0).
\begin{itemize}
\item[(a)] $(H*)^5=(L/C_1)^5$. The eigenvalues of $H*$ are $\mu_i\coloneqq\zeta^iL/C_1$ ($i\in\bbZ/5$), labelled by $\mu_i(0)=\zeta^i$, the classical eigenvalue at $p_i$. The idempotents are
\[
e_i=\frac15\sum_{a=0}^4\Bigl(\prod_{b<a}\kappa^{\mathrm{LR}}_b\Bigr)\mu_i^{-a}H^a,\qquad\Delta_i\coloneqq\frac1{\eta(e_i,e_i)}=\frac{C_0^2\,\zeta^{3i}}{L^2}.
\]
The canonical coordinates along the line are functions $u_i$ with $\partial_Tu_i=\mu_i$, that is $\Dq u_i=\zeta^iL$; they are determined up to additive constants, the constants of $\Dq$ on $\bbQ(\zeta)\fps q[\logq]$ being $\bbQ(\zeta)$ (as in the proof of \cref{lem:set-periods}; for the conifold chart see \cref{lem:set-canon}).
\item[(b)] For every solution $v$ of $z\partial_Tv=H*v$ in a differential $\bbQ\fps q$-algebra in which $z$ is a constant, $C_0v_0$ solves $\PFfive$. In particular, for the column $v=R(z)\tilde e_i\re^{u_i/z}$ of \cref{def:set-rows}, the function $\re^{u_i/z}\,z\,L\,\Phi_i(\zh,u)$ solves $\PFfive$: this is the WKB form of colour $i$.
\end{itemize}
\end{cor}

\begin{proof}
(a) By \cref{prop:set-cyclicLR}~(i), $(H*)^5H^a=\prod_b\kappa^{\mathrm{LR}}_b\,H^a=(L/C_1)^5H^a$, so the characteristic polynomial of $H*$ is $x^5-(L/C_1)^5$, with the pairwise distinct roots $\zeta^iL/C_1$, as $L/C_1\in1+q\bbQ\fps q$. The equation $H*x=\mu_ix$ for $x=\sum_ax_aH^a$ reads $x_a=x_{a-1}\kappa^{\mathrm{LR}}_{a-1}/\mu_i$, consistently around the cycle as $\prod_b\kappa^{\mathrm{LR}}_b=\mu_i^5$. As $*$ is commutative and associative with unit $1$, the eigenlines are ideals spanned by idempotents $e_i$ with $\sum_ie_i=1$; for the coefficient $x_0^{(i)}$ of $H^0$ in $e_i$, the coefficient of $H^a$ in $\sum_ie_i=1$ is $\bigl(\prod_{b<a}\kappa^{\mathrm{LR}}_b\bigr)(C_1/L)^a\sum_ix_0^{(i)}\zeta^{-ia}=\delta_{a0}$, whose unique solution is $x_0^{(i)}=1/5$. By the Frobenius property and $\eta(1,H^a)=5\,\delta_{a3}$, $\eta(e_i,e_i)=\eta(1,e_i)=\kappa^{\mathrm{LR}}_0\kappa^{\mathrm{LR}}_1\kappa^{\mathrm{LR}}_2\,\mu_i^{-3}=L^2/(C_0^2\zeta^{3i})$. Finally $\partial_Tu_i=\mu_i$ and $\Dq=C_1\partial_T$ give $\Dq u_i=\zeta^iL$.

(b) By \cref{prop:set-cyclicLR}~(ii), $v_1=z\partial_Tv_0$, $v_2=z^2(\kappa^{\mathrm{LR}}_1)^{-1}\partial_T^2v_0$, $v_3=z^3\partial_T\bigl((\kappa^{\mathrm{LR}}_1)^{-1}\partial_T^2v_0\bigr)$ and, by \eqref{eq:set-McycLR},
\[
v_4=z^4(\kappa^{\mathrm{LR}}_3)^{-1}\Mcyc^{\mathrm{LR}}(C_0v_0)=z^4\,\PFfour(C_0v_0).
\]
The closing relation $z\partial_Tv_4=\kappa^{\mathrm{LR}}_4v_0$ becomes $z^5C_1^{-1}\Dq\PFfour(C_0v_0)=(C_0/C_1)v_0$, i.e.\ $\PFfive(C_0v_0)=0$ by \cref{lem:set-pf5}. The last statement concerns the rows and is proved after \cref{def:set-rows}.
\end{proof}

\begin{defn}[Points and variables]
\label{def:set-points}
Put $u\coloneqq1-5^5q$ and $s\coloneqq u^{1/5}$, so that $L=u^{-1/5}=s^{-1}$ and $\Dq=q\,\dd/\dd q=-(1-u)\,\dd/\dd u$, and put $\zh\coloneqq z/s^4=L^4z$, the hatted variable. The \emph{large-radius point} is $q=0$ ($u=1$, $L=s=1$, $\zh=z$), the \emph{conifold point} is $u=0$ ($q=5^{-5}$), and the \emph{orbifold point} is $q=\infty$ ($L=0$). Near $q=0$ the functions of the note lie in $\bbQ\fps q[\logq]$; near the conifold point they are Laurent series in $s$, with $\Dq=-\frac{1-s^5}5s^{-4}\frac{\dd}{\dd s}$; the passage is the conifold expansion of \cref{def:set-confexp}.
\end{defn}

\begin{defn}[Frobenius data and the topological field theory]
\label{def:set-frob}
Along the line $t=T(q)H$ let $e_i$ be the idempotents and $\Delta_i=1/\eta(e_i,e_i)$ (\cref{cor:set-cyclicPF5}~(a)), $\eta$ the pairing of \cref{def:set-fq}. The \emph{normalised canonical basis} is $\tilde e_i\coloneqq\Delta_i^{1/2}e_i$, so $\eta(\tilde e_i,\tilde e_j)=\delta_{ij}$, with the square roots $\zeta^{1/2}\coloneqq\zeta^3$ and $\Delta_i^{1/2}\coloneqq\zeta^{4i}\,C_0/L\in\bbQ(\zeta)\fps q$, so that $(\Delta_i^{1/2})^2=\zeta^{8i}C_0^2/L^2=\Delta_i$. At $q=0$, $\Delta_i(0)=W^0_i\coloneqq\prod_{j\neq i}(\lambda_i-\lambda_j)/(5\lambda_i)=\zeta^{3i}$. The \emph{topological field theory of the formal quintic} is the TFT $\TFT$ on $\bbH\otimes\bbQ(\zeta)\fps q$ with $\TFT_{g,n}(\tilde e_i,\dots,\tilde e_i)\coloneqq\Delta_i^{(2g-2+n)/2}=(\zeta^{4i}C_0/L)^{2g-2+n}$ and $\TFT_{g,n}$ vanishing on mixed colours; its unit is $1_{\TFT}\coloneqq\sum_i\Delta_i^{-1/2}\tilde e_i=\sum_ie_i=1$. We also regard $\TFT$ as a TFT on the colour space, with the orthonormal basis $\tilde e_i$ ($i\in\bbZ/5$) and the unit $1_{\TFT}$. By \cref{thm:set-GT}, since $R(0)=1$, $\TFT$ is the degree-zero part of the cohomological field theory along the line.
\end{defn}

% Section 2, part 4 of 8: the R-matrix and its rows; (QRR); the twist sign; h_w rotation; the graph sum; the two forms of the translation; (T).

%-------------------------------------------------------------------------------
\subsection{The \texorpdfstring{$R$}{R}-matrix, its rows and quantum Riemann--Roch}
\label{sec:set-rows}

\begin{defn}[The $R$-matrix and its rows]
\label{def:set-rows}
Let $R(z)=1+\sum_{k\geq1}R_kz^k$, $R_k\in\End(\bbH)\otimes\bbQ(\zeta)\fps q$, be the $R$-matrix of the semisimple cohomological field theory of the formal quintic along the line $t=T(q)H$ \cite[Section~4.1]{GJR18}, in Givental's convention:
\begin{itemize}
\item[(a)] $R$ is symplectic: $R(-z)^T\eta\,R(z)=\eta$, as matrices in the flat basis;
\item[(b)] for every colour $i$ the column $R(z)\tilde e_i\re^{u_i/z}$ solves the quantum differential equation $z\partial_Tv=H*v$, that is, $w\coloneqq R(z)\tilde e_i$ satisfies $z\partial_Tw+\mu_iw=H*w$ (\cref{cor:set-cyclicPF5}~(a));
\item[(c)] in the graph sums of \cref{def:set-graphsum}, $R$ acts with $R^{-1}(\psi)$ on the legs.
\end{itemize}
Its value at $q=0$ is the quoted input (QRR), \cref{thm:set-qrr}, and its graph sum is the cohomological field theory of the formal quintic, the quoted input (T), \cref{thm:set-GT}.

The \emph{rows} are the lowered components $\Qr^{(i)}_a(z)\coloneqq\eta(H^a,R(z)\tilde e_i)\in\bbQ(\zeta)\fps q\fps z$, with the coefficients $\Qr^{(i)}_{a,n}\coloneqq\coef{z^n}\Qr^{(i)}_a$; $\Qr$ is the matrix $(\Qr^{(i)}_a)$, rows indexed by $a$, columns by $i$. With $\zh=L^4z$ we \emph{define} $\Phi_i(\zh,u)\coloneqq C_0\,\Qr^{(i)}_0(z)/(\zeta^iL)$, i.e.\ $\coef{\zh^n}\Phi_i\coloneqq C_0\,\Qr^{(i)}_{0,n}/(\zeta^iL^{4n+1})$, and the coefficients $\mathrm{ph}_n(u)$ by $\Phi_0(\zh,u)=\sum_{n\geq0}\mathrm{ph}_n(u)\,\zh^n$, power series in $q=(1-u)/5^5$; $\mathrm{ph}_0=1$, because $R(0)=1$ and $\eta(1,\tilde e_0)=\Delta_0^{-1/2}=L/C_0$ (\cref{def:set-frob}). The \emph{flat amplitudes} are
\[
A_n\coloneqq\coef{z^n}\bigl(C_0\,\Qr^{(0)}_0\bigr)=L^{4n+1}\,\mathrm{ph}_n(L^{-5}),\qquad\text{so that}\quad\sum_nA_nz^n=L\,\Phi_0(L^4z,L^{-5}).
\]
\end{defn}

\begin{proof}[Proof of \cref{prop:set-cyclicLR}~(iii) and of the end of \cref{cor:set-cyclicPF5}~(b)]
For $w\coloneqq R(z)\tilde e_i$, \cref{def:set-rows}~(b) and the argument of \cref{prop:set-cyclicLR}~(ii) give $z\partial_Tw_a+\mu_iw_a=\eta(H*H^a,w)=\kappa^{\mathrm{LR}}_aw_{a+1}$, $w_a=\Qr^{(i)}_a$; multiplying by $C_1$, with $C_1\mu_i=\zeta^iL$, gives $(z\Dq+\zeta^iL)\Qr^{(i)}_a=C_1\kappa^{\mathrm{LR}}_a\Qr^{(i)}_{a+1}$. Dividing by $\zeta^i$ gives the relations for $X_a$ ($C_1\kappa^{\mathrm{LR}}_4=C_0$), and $X_0=L\Phi_i/C_0$ is the definition of $\Phi_i$. For the WKB form, $v=R(z)\tilde e_i\re^{u_i/z}$, with the formal exponential $\Dq\re^{u_i/z}=z^{-1}(\Dq u_i)\re^{u_i/z}$, solves $z\partial_Tv=H*v$, so $C_0v_0=\zeta^i\re^{u_i/z}L\Phi_i$ solves $\PFfive$ by \cref{cor:set-cyclicPF5}~(b); multiply by the constant $z\zeta^{-i}$.
\end{proof}

The colour covariance $\Qr^{(i)}_{a,n}=\zeta^{i(a+1-n)}\Qr^{(0)}_{a,n}$ is \cref{lem:gau-colrows}; the equation of $\Phi_0$ is \cref{lem:gau-HPF}.

\begin{thm}[The $R$-matrix at $q=0$, (QRR); quoted]
\label{thm:set-qrr}
For a finite multiset $\mathbf w$ of non-zero numbers put
\begin{equation}
\label{eq:set-hw}
\hw(z)\coloneqq\exp\Bigl(-\sum_{k\geq1}\bk_k\,z^{2k-1}\sum_{a\in\mathbf w}a^{1-2k}\Bigr),\qquad\bk_k\coloneqq\frac{B_{2k}}{2k(2k-1)},
\end{equation}
and $\bw_i\coloneqq\{\zeta^i-\zeta^j\}_{j\neq i}\cup\{-5\zeta^i\}$. In Givental's convention (\cref{def:set-rows}), $R(q=0)=\diag(h_{\bw_0},\dots,h_{\bw_4})$ in the basis $\tilde e_i(q=0)$: $R(0,z)\tilde e_i(0)=h_{\bw_i}(z)\,\tilde e_i(0)$, where $R(0,z)$ is the part of order $q^0$ of $R(z)$, to all orders in $z$.
\end{thm}

This is \cite[proof of Lemma~5.2, (17)]{GJR18} and \cite[proof of Lemma~4.6, (19)--(20)]{GJR17} at $t=0$; both refer to \cite[Section~2.2, Theorem~1 and the Remark after it]{Giv01} and \cite[Theorem~1, Corollary~1]{MR2276766}. It is used only in the following form, for every colour $i$:
\begin{equation}
\label{eq:set-qrr}
\Phi_i(z,1)=h_{\bw_i}(z);
\end{equation}
in particular, with $h\coloneqq h_{\bw_0}$ and $h_n\coloneqq\coef{z^n}h$, $A_n(q=0)=\mathrm{ph}_n(1)=h_n$.

\begin{proof}[Proof of \eqref{eq:set-qrr}]
At $q=0$, $u=1$, $L=1$, $\zh=z$ and $C_0=1$. For every $q$, $\eta(1,\tilde e_i)=\Delta_i^{1/2}\eta(e_i,e_i)=\Delta_i^{-1/2}=\zeta^iL/C_0$ (\cref{cor:set-cyclicPF5}~(a), \cref{def:set-frob}). By (QRR) the part of order $q^0$ of $\Qr^{(i)}_0$ is $h_{\bw_i}(z)\,\eta(1,\tilde e_i)|_{q=0}=\zeta^ih_{\bw_i}(z)$, so $\Phi_i(z,1)=h_{\bw_i}(z)$.
\end{proof}

\begin{rmk}[The twist sign]
\label{rmk:set-qrrsign}
The weight $-5\zeta^i$ in $\bw_i$ is minus the fibre weight $5\zeta^i$ of $\cO(5)$ at $p_i$, the sign of an Euler-class twist: the Chern root $5\zeta^i$ contributes $\exp\bigl(+\sum_m\bk_m(z/5\zeta^i)^{2m-1}\bigr)=\exp\bigl(-\sum_m\bk_mz^{2m-1}(-5\zeta^i)^{1-2m}\bigr)$. The same sign appears in \cite[(17)]{GJR18}, \cite[Corollary~1]{MR2276766} and \cite[(26)]{LP20}. Pairing $j$ with $5-j$ gives $\frac1{1-\zeta^j}+\frac1{1-\zeta^{-j}}=1$, so $\sum_{a\in\bw_0}a^{-1}=2-\frac15=\frac95$ and $h_1=-\frac1{12}\cdot\frac95=-\frac3{20}$, whereas the weight $+5$ would give $-\frac{11}{60}$. That $\mathrm{ph}_1=-\frac3{20}$ is forced, independently of (QRR), is part of \cref{lem:gau-HPF}.
\end{rmk}

\begin{lem}[Rotation of the colours]
\label{lem:set-hrot}
$h_{\bw_i}(z)=h_{\bw_0}(\zeta^{-i}z)$ for every $i$.
\end{lem}

\begin{proof}
As multisets $\bw_i=\zeta^i\bw_0$, so $\sum_{a\in\bw_i}a^{1-2k}z^{2k-1}=\sum_{a\in\bw_0}a^{1-2k}(\zeta^{-i}z)^{2k-1}$ in \eqref{eq:set-hw}.
\end{proof}

%-------------------------------------------------------------------------------
\subsection{The graph sum}
\label{sec:set-graph}

\begin{defn}[The Givental--Teleman graph sum]
\label{def:set-graphsum}
Let $\Omega$ be a cohomological field theory with unit $1$ \cite[Def.~0.3]{PPZ15} on a free module $V$ of finite rank with pairing $\eta$ over a commutative $\bbQ$-algebra, and $R(z)=1+O(z)\in\End(V)\fps z$ symplectic, $R(-z)^*R(z)=1$ with ${}^*$ the $\eta$-adjoint \cite[(5)]{PPZ15}. The ancestor correlators are $\langle\psi^{k_1}a_1\cdots\psi^{k_n}a_n\rangle_h\coloneqq\int_{\overline{\mathcal M}_{h,n}}\Omega_{h,n}(a_1,\dots,a_n)\prod_j\psi_j^{k_j}$; for $q\in V\fps z$ (a point of the loop space, not the parameter $q$) the dilaton-shifted potential is $\overline F^\Omega_h(q)\coloneqq\sum_n\frac1{n!}\langle(q+z1)^n\rangle_h$, with $z^ka\mapsto\psi^ka$ (base point $q=-z1$), and $\dd^n\overline F^\Omega_h(q)$ is the symmetric form $(x_1,\dots,x_n)\mapsto\sum_m\frac1{m!}\langle x_1\cdots x_n(q+z1)^m\rangle_h$ on $V\fps z$. The \emph{edge} and the \emph{leaf point} of $R$ are
\begin{gather*}
E_R(\psi,\psi')\coloneqq\frac{\eta^{-1}-\bigl(R^{-1}(\psi)\otimes R^{-1}(\psi')\bigr)\eta^{-1}}{\psi+\psi'}\in(V\otimes V)\fps{\psi,\psi'},\\
q_R(z)\coloneqq R^{-1}(z)(-z1)=-z1+\ttr_R(z),
\end{gather*}
where $\eta^{-1}\in V\otimes V$ is the inverse bivector, $E_R$ is a power series by the symplectic condition \cite[Rem.~2.1, (6)]{PPZ15}, and $\ttr_R(z)\coloneqq z1-zR^{-1}(z)1\in z^2V\fps z$ is the translation of \cite[Def.~2.13]{PPZ15}. For a stable graph $\Gamma$ without legs, with vertices $v$ of genus $g_v$ and valence $n_v$ and of total genus $g=\sum_vg_v+h^1(\Gamma)$, the contraction
\[
\mathrm{Cont}_\Gamma[E,q;\Omega]\coloneqq|\Aut\Gamma|^{-1}\Bigl\langle\prod_v\dd^{n_v}\overline F^\Omega_{g_v}(q),\ \prod_eE\Bigr\rangle
\]
contracts, at each vertex, the $n_v$-th derivative of $\overline F^\Omega_{g_v}$ at $q$, one slot for each half-edge, with one copy of $E(\psi,\psi')$ for each edge, its two slots at the two half-edges; sums over graphs carry no further factor. For $q=q_R$ the contraction is finite, as $\ttr_R\in z^2V\fps z$. Let $R.\Omega$ be the unit-preserving action of \cite[Defs.~2.2, 2.5, 2.13]{PPZ15}: legs $R^{-1}(\psi)$, edges $E_R$, translation $\ttr_R$, as in Givental and Teleman \cite[p.~20]{PPZ15}. There the translation points are pushed forward, whereas the vertex tensors above insert them; for $g\geq2$ the two forms agree (\cref{lem:set-pushins}~(iii)). The sum $\sum_\Gamma\mathrm{Cont}_\Gamma[E_R,q_R;\Omega]$ over the leg-free stable graphs of genus $g$ is the \emph{graph sum of genus $g$}. The definitions and propositions of \cite{PPZ15} used in the note are identities between finite sums, polynomial in the coefficients; stated there over $\bbQ$, they hold verbatim over every commutative $\bbQ$-algebra.

\emph{Colours.} When $\Omega=\TFT$ (\cref{def:set-frob}), which vanishes on mixed colours, a \emph{coloured graph} is a stable graph with a colour at each vertex, and the term of $\Gamma$ is the sum over its colourings. For a constant isometry $A\colon(W,\eta_W)\to(V,\eta)$ and $A_*\Omega\coloneqq\Omega\circ(A^{-1})^{\otimes n}$, $A_*(R.\Omega)=(ARA^{-1}).(A_*\Omega)$, immediate from \cite[Defs.~2.2, 2.13]{PPZ15}; it gives the graph sum in the colour basis (\cref{sec:sgn}).
\end{defn}

% Section 2, part 5 of 8: the translation (push-forward and insertion), (T); the Frobenius basis at the conifold point; conifold frames; the differential ring and the conifold expansion; tau; the canonical coordinates.

\begin{lem}[The translation: push-forward and insertion]
\label{lem:set-pushins}
Let $\Omega$ be a cohomological field theory \cite[Def.~0.3]{PPZ15} on a free module $V$ of finite rank with pairing $\eta$ over a commutative $\bbQ$-algebra, with the correlators $\langle\cdots\rangle_h$ of \cref{def:set-graphsum}. Let $2h-2+n>0$, $m\geq0$, and let $p_m\colon\overline{\mathcal M}_{h,n+m}\to\overline{\mathcal M}_{h,n}$ forget the last $m$ points.
\begin{enumerate}
\item For $k_1,\dots,k_n\geq0$ and $c_1,\dots,c_m\geq1$, in $H^*(\overline{\mathcal M}_{h,n+m})$,
\[
\prod_{j=1}^n\psi_j^{k_j}\prod_{s=1}^m\psi_{n+s}^{c_s}=\prod_{j=1}^n(p_m^*\psi_j)^{k_j}\prod_{s=1}^m\psi_{n+s}^{c_s}.
\]
\item Let $\ttr\in z^2V\fps z$ and let $\ttr\Omega$ be the translation of $\Omega$ by $\ttr$ \cite[Def.~2.5]{PPZ15}. For $a_1,\dots,a_n\in V$ and $k_1,\dots,k_n\geq0$,
\[
(\ttr\Omega)_{h,n}(a_1,\dots,a_n)\prod_{j=1}^n\psi_j^{k_j}=\sum_{m\geq0}\frac1{m!}\,p_{m*}\Bigl(\Omega_{h,n+m}\bigl(a_1,\dots,a_n,\ttr(\psi_{n+1}),\dots,\ttr(\psi_{n+m})\bigr)\prod_{j=1}^n\psi_j^{k_j}\Bigr),
\]
where on the right $\psi_j$ is the class of $\overline{\mathcal M}_{h,n+m}$; consequently
\[
\int_{\overline{\mathcal M}_{h,n}}(\ttr\Omega)_{h,n}(a_1,\dots,a_n)\prod_{j=1}^n\psi_j^{k_j}=\sum_{m\geq0}\frac1{m!}\bigl\langle\psi^{k_1}a_1\cdots\psi^{k_n}a_n\,\ttr^m\bigr\rangle_h ,
\]
$\ttr$ being inserted at $m$ further points through $z^ca\mapsto\psi^ca$. The sums are finite. If $\Omega$ has the unit $1$, the right side is $\dd^n\overline F^\Omega_h(q)(\psi^{k_1}a_1,\dots,\psi^{k_n}a_n)$ at $q=-z1+\ttr$.
\item Let $\Omega$ have the unit $1$ and let $R(z)=1+O(z)$ be symplectic, as in \cref{def:set-graphsum}. For $g\geq2$ and every leg-free stable graph $\Gamma$ of genus $g$, $\mathrm{Cont}_\Gamma[E_R,q_R;\Omega]$ is the term of $\Gamma$ in $\int_{\overline{\mathcal M}_g}(R.\Omega)_{g,0}$, that is, $|\Aut\Gamma|^{-1}$ times the integral over $\overline{\mathcal M}_g$ of the contribution of $\Gamma$ in \cite[Def.~2.2]{PPZ15} to the action of $R$ on $\ttr_R\Omega$, which is $R.\Omega$ \cite[Def.~2.13]{PPZ15}. Consequently $\int_{\overline{\mathcal M}_g}(R.\Omega)_{g,0}$ is the graph sum of genus $g$. By \cite[Prop.~2.12]{PPZ15}, $R.\Omega$ is again a cohomological field theory with the unit $1$, to which (iii) applies.
\end{enumerate}
\end{lem}

\begin{proof}
(i) For $2h-2+N>0$ let $\pi\colon\overline{\mathcal M}_{h,N+1}\to\overline{\mathcal M}_{h,N}$ forget the last point and, for $j\leq N$, let $\delta_j$ be the divisor of the curves with a genus-$0$ component carrying only the markings $j$ and $N+1$ (local letters). The formula \cite[(8)]{PPZ15}, $\pi^*(\psi_j^d)=\psi_j^d-\delta_j\,\pi^*(\psi_j^{d-1})$, gives $\pi^*\psi_j=\psi_j-\delta_j$ for $d=1$, and for $d=2$, compared with $(\pi^*\psi_j)^2=(\psi_j-\delta_j)^2$, $\psi_j\delta_j=0$. Exchanging the markings $j$ and $N+1$ preserves $\delta_j$ and exchanges $\psi_j$ and $\psi_{N+1}$; so $\psi_{N+1}\delta_j=0$, i.e.\ $\psi_j\psi_{N+1}=\pi^*(\psi_j)\,\psi_{N+1}$. Replacing the factors $\psi_i$ one at a time, for $k_i\geq0$ and $c\geq1$,
\[
\textstyle\prod_{i=1}^N\psi_i^{k_i}\,\psi_{N+1}^c=\prod_{i=1}^N(\pi^*\psi_i)^{k_i}\,\psi_{N+1}^c,\tag{$*$}
\]
which for $N=n$ is (i) with $m=1$. For $m\geq2$ write $p_m=p_{m-1}\circ\pi$, $\pi$ forgetting the last point of $\overline{\mathcal M}_{h,n+m}$, and let $Y\coloneqq\prod_{s<m}\psi_{n+s}^{c_s}$, $X\coloneqq\prod_j\psi_j^{k_j}\,Y$ on $\overline{\mathcal M}_{h,n+m-1}$ (local letters). By $(*)$ with $N=n+m-1$ the left side of (i) is $\pi^*(X)\,\psi_{n+m}^{c_m}$; by (i) for $m-1$, $\pi^*(X)=\prod_j(p_m^*\psi_j)^{k_j}\,\pi^*(Y)$; and by $(*)$ again $\pi^*(Y)\,\psi_{n+m}^{c_m}=\prod_{s\leq m}\psi_{n+s}^{c_s}$ on $\overline{\mathcal M}_{h,n+m}$.

(ii) By \cite[Def.~2.5]{PPZ15},
\[
(\ttr\Omega)_{h,n}(a_1,\dots,a_n)=\sum_m\frac1{m!}\,p_{m*}Z_m,\qquad Z_m\coloneqq\Omega_{h,n+m}(a_1,\dots,a_n,\ttr(\psi_{n+1}),\dots,\ttr(\psi_{n+m}))
\]
(a local letter), an argument $\ttr(\psi_{n+s})$ standing for $\sum_{c\geq2}\psi_{n+s}^c$ times the class with the argument $\ttr_c\coloneqq\coef{z^c}\ttr$. By the projection formula, $p_{m*}(Z_m)\prod_j\psi_j^{k_j}=p_{m*}\bigl(Z_m\prod_j(p_m^*\psi_j)^{k_j}\bigr)$, and as every term of $Z_m$ carries $\prod_s\psi_{n+s}^{c_s}$ with $c_s\geq2$, (i) replaces $p_m^*\psi_j$ by $\psi_j$. This is the first identity; the second follows from $\int p_{m*}W=\int W$. A term with $m$ translation points has degree at least $2m$ on a space of dimension $3h-3+n+m$, so the sums are finite. The last sentence is the definition of $\dd^n\overline F^\Omega_h$, as $q+z1=\ttr$.

(iii) By \cite[Defs.~2.2, 2.13]{PPZ15}, $(R.\Omega)_{g,0}=\sum_\Gamma|\Aut\Gamma|^{-1}C_\Gamma$ (a local letter) over the leg-free stable graphs of genus $g$, where $C_\Gamma$ is the push-forward along the gluing map $\prod_v\overline{\mathcal M}_{g_v,n_v}\to\overline{\mathcal M}_g$ \cite[\S0.2, (1)]{PPZ15} of the contraction of the classes $(\ttr_R\Omega)_{g_v,n_v}$ at the vertices with the edges $E_R(\psi',\psi'')$ in the classes $\psi'$, $\psi''$ of the two half-edges \cite[Def.~2.2~(iii)]{PPZ15}, formed with $R^{-1}$ \cite[p.~20]{PPZ15}. The push-forward preserves integrals. Expanding each edge into finitely many relevant monomials $\psi'^{\,a}\psi''^{\,b}$, the integral over the product is the contraction of the forms $(x_1,\dots,x_{n_v})\mapsto\int_{\overline{\mathcal M}_{g_v,n_v}}(\ttr_R\Omega)_{g_v,n_v}(x_1,\dots,x_{n_v})\prod_l\psi_l^{k_l}$, which by (ii), as $q_R+z1=\ttr_R$, are $\dd^{n_v}\overline F^\Omega_{g_v}(q_R)(\psi^{k_1}x_1,\dots,\psi^{k_{n_v}}x_{n_v})$. Summing over the monomials gives $\mathrm{Cont}_\Gamma[E_R,q_R;\Omega]$, and summing over $\Gamma$ the graph sum.
\end{proof}

\begin{thm}[The reconstruction (T); quoted]
\label{thm:set-GT}
The cohomological field theory of the formal quintic along the line, i.e.\ the classes $\Omega_{g,n}$ of \cref{def:set-fq} with $\Qn=\re^{T(q)}$ (which, by the divisor equation, is the theory at the point $t=T(q)H$ with $\Qn=1$), equals $R.\TFT$, with $R$ of \cref{def:set-rows} and $\TFT$ of \cref{def:set-frob} \cite[Theorem~(2) in (1.15)]{MR2917177}, \cite[Section~4.1, Theorem~4.1]{GJR18}. In particular, for $g\geq2$, since $F_g=\int_{\overline{\mathcal M}_g}\Omega_{g,0}$ (\cref{def:set-fq}), \cref{lem:set-pushins}~(iii) gives
\[
F_g=\sum_\Gamma\mathrm{Cont}_\Gamma[E_R,q_R;\TFT]
\]
as power series in $q$, the sum running over the leg-free stable graphs of genus $g$.
\end{thm}

%-------------------------------------------------------------------------------
\subsection{Conifold frames and the conifold expansion}
\label{sec:set-frames}

In this subsection functions of $u$ are formal series in $u$ (or in $s=u^{1/5}$), $\Dq=-(1-u)\,\dd/\dd u$, and $\PFfour$ is taken in its $u$-form (\cref{def:set-genus0}).

\begin{lem}[The Frobenius basis at the conifold point]
\label{lem:set-frobbasis}
Let $\theta\coloneqq u\,\dd/\dd u$, so that $\Dq=(1-u^{-1})\theta$ on $\bbQ\lser u$, and let $\ell$ be an indeterminate with $\Dq\ell=1-u^{-1}$, i.e.\ $\theta\ell=1$, a formal $\log u$ (the indeterminate of \cref{def:gau-conc}).
\begin{enumerate}
\item $\PFfour=\sum_{j=-3}^{1}u^jP_j(\theta)$ with polynomials $P_j$, where
\[
P_{-3}(x)=x(x-1)^2(x-2),\qquad P_{-2}(x)=-x(x-1)\bigl(4x^2-6x+\tfrac{17}5\bigr),\qquad P_{-1}(0)=0,\quad P_{-1}'(0)=-\tfrac75 .
\]
Hence $\PFfour$ maps $\bbQ\fps u$ into itself, its local exponents at $u=0$ are $0,1,1,2$, and a power-series solution is determined by its coefficients of $u^0,u^1,u^2$; $f_k$ ($k=0,1,2$) denotes the solution with $f_k=u^k+O(u^3)$.
\item For $g=\sum_nb_nu^n\in\bbQ\fps u$, $\PFfour(g\ell)=\ell\,\PFfour g+r(g)$ with $r(g)\coloneqq\sum_nb_n\sum_jP_j'(n)u^{n+j}\in u^{-3}\bbQ\fps u$. If $\PFfour g=0$, then $g\ell+h$ with $h\in\bbQ\fps u$ solves $\PFfour$ for some $h$ if and only if $\coef{u^m}r(g)=0$ for $m=-3,-2,-1$.
\item On $\operatorname{span}(f_0,f_1,f_2)$ the obstruction $(\coef{u^{-3}}r,\coef{u^{-2}}r,\coef{u^{-1}}r)$ has the matrix, in the columns $f_0,f_1,f_2$,
\[
\begin{pmatrix}-2&0&0\\ \frac{17}5&0&0\\ -\frac75&-\frac75&2\end{pmatrix},
\]
with kernel $\bbQ\,(0,1,\frac7{10})$. Hence, up to a non-zero factor, $\varpi_1\coloneqq f_1+\gamma f_2$, $\gamma\coloneqq\frac7{10}$, is the only power-series solution $g$ for which a solution $g\ell+h$, $h\in\bbQ\fps u$, exists; $h$ is unique up to adding power-series solutions, and the logarithmic solution with $\coef{u^0}h=\coef{u^1}h=\coef{u^2}h=0$ is $\varpi_1\ell+\rho_3$, $\rho_3\in u^3\bbQ\fps u$. The \emph{vanishing period} $\varpi_1=u+\frac7{10}u^2+O(u^3)$ is normalised by $\varpi_1=u+O(u^2)$.
\end{enumerate}
\end{lem}

\begin{proof}
(i) As $\theta u^m=u^m(\theta+m)$, every $\sum_kc_k\Dq^k$, $c_k\in\bbQ[u,u^{-1}]$, is uniquely $\sum_ju^jP_j(\theta)$. Moving the factors $u^{-1}$ to the left, the part of order $u^{-k}$ of $\Dq^k$ is $(-1)^ku^{-k}\theta(\theta-1)\cdots(\theta-k+1)$, and its part of order $u^{1-k}$ is $(-1)^{k-1}u^{1-k}\theta(\theta-1)\cdots(\theta-k+2)\bigl(k\theta-\binom k2\bigr)$ (the word with $\theta$ in position $p+1$ contributes $\theta\cdots(\theta-k+2)(\theta-k+1+p)$, $p=0,\dots,k-1$). In the $u$-form of $\PFfour$ the terms of order $u^{-3}$ come from $u\Dq^4$ and $-2\Dq^3$, those of order $u^{-2}$ from $u\Dq^4$, $-2(1-u)\Dq^3$ and $-\frac75\Dq^2$:
\begin{align*}
P_{-3}(\theta)&=\theta(\theta-1)(\theta-2)(\theta-3)+2\theta(\theta-1)(\theta-2)=\theta(\theta-1)^2(\theta-2),\\
P_{-2}(\theta)&=-\theta(\theta-1)(\theta-2)(4\theta-6)-2\theta(\theta-1)(\theta-2)-6\theta(\theta-1)^2-\tfrac75\theta(\theta-1)\\
&=-\theta(\theta-1)\bigl(4\theta^2-6\theta+\tfrac{17}5\bigr).
\end{align*}
Every $\Dq^k$, $k\geq1$, ends with $\theta$, and the constant term $-\frac{24}{625}(1-u)$ contributes to $P_0$, $P_1$ only; so $P_j(0)=0$ for $j\leq-1$. Since $\PFfour u^n=\sum_jP_j(n)u^{n+j}$ and $P_j(n)=0$ whenever $n\geq0>n+j$ (the roots $0,1,2$ of $P_{-3}$, $0,1$ of $P_{-2}$, $0$ of $P_{-1}$), $\PFfour$ preserves $\bbQ\fps u$; the coefficient of $u^N$ in $\PFfour g$ is $P_{-3}(N+3)b_{N+3}$, with $P_{-3}(N+3)=(N+3)(N+2)^2(N+1)\neq0$, plus a combination of $b_{N-1},\dots,b_{N+2}$, so $b_0,b_1,b_2$ are free and determine $g$.

(ii) As $\theta(u^n\ell)=u^n(n\ell+1)$, $P(\theta)(u^n\ell)=u^n(P(n)\ell+P'(n))$ for every polynomial $P$; summing gives $\PFfour(g\ell)=\ell\,\PFfour g+r(g)$, and for $n\geq3$ all exponents $n+j$ are $\geq0$. If $\PFfour g=0$, $\PFfour(g\ell+h)=r(g)+\PFfour h$ with $\PFfour h\in\bbQ\fps u$; so a solution requires $\coef{u^m}r(g)=0$ for $m<0$, and conversely $\PFfour h=-r(g)$ is then solved by the recursion of (i).

(iii) $\coef{u^m}r(g)=\sum_jP_j'(m-j)\,b_{m-j}$ over $m-j\geq0$: for $m=-3$, $P_{-3}'(0)b_0=-2b_0$; for $m=-2$, $P_{-3}'(1)b_1+P_{-2}'(0)b_0=\frac{17}5b_0$ ($1$ is a double root of $P_{-3}$); for $m=-1$, $2b_2-\frac75b_1+P_{-1}'(0)b_0$. And $P_{-1}'(0)=\coef{u^{-1}}\PFfour(\ell)$: from $\Dq\ell=1-u^{-1}$, $\Dq^2\ell=u^{-1}-u^{-2}$, $\Dq^3\ell=-u^{-1}+3u^{-2}-2u^{-3}$, $\Dq^4\ell=u^{-1}-7u^{-2}+12u^{-3}-6u^{-4}$, this coefficient is $-7+(2+6)-\frac75-\frac75+\frac25=-\frac75$. The kernel consists of the $(b_0,b_1,b_2)$ with $b_0=0$, $2b_2=\frac75b_1$; the rest follows from (ii), as two logarithmic solutions with the same $g$ differ by a power-series solution.
\end{proof}

% Section 2, part 6 of 8: conifold frames; the differential ring and the conifold expansion; tau; the canonical coordinates; the model, its rows and the T-system.

\begin{defn}[Conifold frames]
\label{def:set-frames}
A \emph{conifold frame} is given by $(\alpha,\beta)\in\bbQ^2$: $\varpi_0\coloneqq f_0+\alpha\varpi_1+\beta f_2$ and $\tc\coloneqq\varpi_1/\varpi_0$. Then $\varpi_0=1+\alpha u+O(u^2)$ and $\tc=u+(\frac7{10}-\alpha)u^2+O(u^3)$, so $\varpi_0(0)=1$, $\tc'(0)=1$ and $\Dq\tc(0)=-1$; in the frame $(0,0)$, $\tc=u+\frac7{10}u^2+\frac{41}{75}u^3+O(u^4)$. These are \emph{the conifold frames} of the note. In them the normalisation (N0) of \cref{lem:gau-diag} holds; it uses only $\varpi_0(0)=1$, $\tc'(0)=1$, $\varpi_1=u+O(u^2)$, and $v(u)=u+O(u^2)$ (\cref{def:set-tau}).
\end{defn}

\begin{defn}[The differential ring of the periods and the conifold expansion]
\label{def:set-confexp}
Let
\[
\Nring\coloneqq\bbQ[L^{\pm1},\logq,c_0,c_1,c_2,c_3,i_0,i_1,i_2][c_0^{-1},\mathsf W^{-1}],\qquad\mathsf W\coloneqq c_0^2+c_0i_1-c_1i_0,
\]
a localised polynomial ring, with the derivation $\Dq$ fixed on the generators: $\Dq L\coloneqq(L^6-L)/5$, $\Dq\logq\coloneqq1$, $\Dq c_j\coloneqq c_{j+1}$ ($j\leq2$), $\Dq c_3\coloneqq(L^5-1)\bigl(2c_3+\frac75c_2+\frac25c_1+\frac{24}{625}c_0\bigr)$, which is $\PFfour c_0=0$ with $u=L^{-5}$, $\Dq i_j\coloneqq i_{j+1}$ ($j\leq1$), and $\Dq i_2\coloneqq i_3$, the element of $\Nring$ determined by the concomitant relation $\Conc(c_0,I_1)=0$, where $I_1\coloneqq c_0\logq+i_0$, $\Dq^jI_1=\logq\,c_j+jc_{j-1}+i_j$ for $j\leq3$, and $\Conc$ is the concomitant of $\PFfour$ (\cref{def:gau-conc}): the terms with $\logq$ cancel as $\Conc(c_0,c_0)=0$, and the relation is linear in $i_3$ with the coefficient $-uc_0=-L^{-5}c_0$, a unit. Put $C_0\coloneqq c_0$, $T\coloneqq I_1/C_0$ and $C_1\coloneqq\Dq T=\mathsf W/c_0^2$. The \emph{evaluation} $\ev\colon\Nring\to\bbQ\fps q[\logq]$, $c_j\mapsto\Dq^jC_0$, $i_j\mapsto\Dq^j\tilde I_1$, $\logq\mapsto\logq$, $L\mapsto L$, is a homomorphism of differential rings, with $\ev(\mathsf W)=C_0^2C_1$: it respects $\Dq c_3$ as $\PFfour C_0=0$, and $\Dq i_2$ as $\Conc(C_0,I_1)=0$; indeed $\PFfour C_0=\PFfour I_1=0$ (\cref{lem:set-periods}), so \cref{lem:gau-conc}~(i), with \cref{lem:gau-adj}, gives $\Dq\Conc(C_0,I_1)=0$, and this constant vanishes, as its part of order $q^0$ involves only the leading coefficient of $\PFfour$ and the jets $(1,0,0,0)$, $(\logq,1,0,0)$ of $C_0$, $I_1$.

The \emph{conifold expansion} in the frame $(\alpha,\beta)$ is the ring homomorphism $\Cexp_{\alpha,\beta}\colon\Nring\to\bbQ\lser s$ with
\[
c_j\mapsto\Dq^j\varpi_0,\qquad i_j\mapsto\Dq^j\varpi_1-j\,\Dq^{j-1}\varpi_0+\Xi\,\Dq^j\varpi_0,\qquad\logq\mapsto\log(1-u)=-\Xi,\qquad L\mapsto s^{-1},
\]
where $\Xi\coloneqq-\log(1-u)=\sum_{n\geq1}u^n/n$; it is well defined because $c_0\mapsto\varpi_0$ and $\mathsf W\mapsto\varpi_0\Dq\varpi_1-\varpi_1\Dq\varpi_0=\varpi_0^2\Dq\tc$ are units ($\varpi_0(0)=1$, $\Dq\tc(0)=-1$). It maps $C_0$, $I_1$, $T$, $C_1$ to $\varpi_0$, $\varpi_1$, $\tc$, $\Dq\tc$. Replacing the primitive $-\Xi$ of $1$ by $-\Xi+k$, $k\in\bbQ$, composes $\Cexp_{\alpha,\beta}$ with the automorphism $i_j\mapsto i_j-kc_j$, $\logq\mapsto\logq+k$ of $\Nring$, which commutes with $\Dq$ and fixes $c_0$, $I_1$, $\mathsf W$ and $L$, hence $C_0$, $C_1$, $\Delta_i^{1/2}=\zeta^{4i}C_0/L$, the rows (\cref{cor:gau-rowsNN}) and the genus-$g$ sum $F_g^{\Nring}$ of \cref{def:mai-Fcf}. That $\Cexp_{\alpha,\beta}$ commutes with $\Dq$ is \cref{lem:gau-confexp}, and that the rows lie in $\Nring\otimes\bbQ(\zeta)$ is \cref{cor:gau-rowsNN}. An element of $\Nring\otimes\bbQ(\zeta)$, or a series over it, is \emph{regular at the conifold point} if its conifold expansion lies in $\bbQ\fps u$, coefficientwise. The comparison with the conifold expansion of \cite{CGYZ26} is \cref{rmk:mai-constant}.
\end{defn}

\begin{defn}[The natural variable]
\label{def:set-tau}
Put
\[
\tau(u)\coloneqq\frac1{{}_2F_1(1,\frac45;\frac95;u)}=\Bigl(\sum_{n\geq0}\frac{4u^n}{5n+4}\Bigr)^{-1}\in1+u\bbQ\fps u,\qquad v(u)\coloneqq u\,\tau(u)^{-5/4}\in u+u^2\bbQ\fps u,
\]
the power being the binomial series, and $\varsigma\coloneqq v^{1/5}=s\,\tau^{-1/4}$, so that $z/\varsigma^4=\tau(u)\zh$; $\tau=1-\frac49u+O(u^2)$, and the letter $\tau$ denotes only this series.
\end{defn}

\begin{lem}[The canonical coordinates at the conifold point]
\label{lem:set-canon}
In $\bbQ(\zeta)\lser s$, with $\Dq=-(1-u)\,\dd/\dd u$ and $L=s^{-1}$, the solutions $y$ of $\Dq y=\zeta^iL$ are the series $y=c-\frac54\zeta^iu^{4/5}/\tau(u)$, $c\in\bbQ(\zeta)$, where $u^{4/5}=s^4$. So in the conifold chart the canonical coordinates of \cref{cor:set-cyclicPF5}~(a), which the note uses only through the equation $\Dq u_i=\zeta^iL$ and which are determined by it up to additive constants, are the series $u_i=u_i^{\mathrm c}-\frac54\zeta^iu^{4/5}/\tau(u)$, with constants $u_i^{\mathrm c}$, their values at $u=0$; and
\[
\tau(u)\,\zh=-\frac54\,\frac{\zeta^iz}{u_i-u_i^{\mathrm c}},\qquad u_i-u_i^{\mathrm c}=\zeta^i\bigl(u_0-u_0^{\mathrm c}\bigr).
\]
\end{lem}

\begin{proof}
Termwise, $\int_0^ut^{-1/5}(1-t)^{-1}\dd t=\sum_n\frac{u^{n+4/5}}{n+4/5}=\frac54u^{4/5}/\tau(u)$, so $\Dq\bigl(-\frac54\zeta^iu^{4/5}/\tau\bigr)=\zeta^iu^{-1/5}=\zeta^iL$; the constants of $\Dq$ on $\bbQ(\zeta)\lser s$ are $\bbQ(\zeta)$, since $\Dq f=0$ means $\dd f/\dd s=0$. With $\Dq u_i=\zeta^iL$ (\cref{cor:set-cyclicPF5}~(a)) and $\zh=z/u^{4/5}$ the rest follows.
\end{proof}

The constants $u_i^{\mathrm c}$ play no role: the exponentials $\re^{u_i/z}$ are stripped from the formal solutions and never absorbed into a normalisation series (\cref{thm:gau-R0}).

%-------------------------------------------------------------------------------
\subsection{The model}
\label{sec:set-model}

\begin{defn}[The model]
\label{def:set-model}
The \emph{model} is the Landau--Ginzburg potential $\Wv(p)\coloneqq p^4/4+v/p$, with the parameter $v$ and the primitive form $\dd p$. Its formal integrals $I_m(v,z)\coloneqq\int p^m\re^{-\Wv(p)/z}\dd p$ ($m\in\bbZ$) are symbols subject to the two relations
\[
\text{(Mod1)}\quad z\,\partial_vI_m=-I_{m-1},\qquad\text{(Mod2)}\quad mz\,I_{m-1}-I_{m+3}+v\,I_{m-2}=0,
\]
which express the differentiation under the integral sign ($\partial_v\Wv=p^{-1}$) and the vanishing of the integral of $\dd(p^m\re^{-\Wv/z})=\bigl(mp^{m-1}-(p^{m+3}-vp^{m-2})/z\bigr)\re^{-\Wv/z}\dd p$; they are realised by residues at $p=0$ (\cref{def:gau-van}) and by the formal stationary phase at a critical point (\cref{lem:con-S2}). The charges and the signs are $\nu\coloneqq(0,-1,3,2,1)$ and $\varepsilon\coloneqq(1,-1,-1,1,1)$, with indices in $\bbZ/5$; thus $\nu_a\equiv-a\bmod5$ and $\nu_a+\nu_{3-a}=2$. The \emph{rows} of the model are $Y_b\coloneqq\varepsilon_bI_{\nu_b}$, i.e.\ $Y=(I_0,-I_{-1},-I_3,I_2,I_1)$, and the dictionary of classes is $H^b\leftrightarrow\varepsilon_b\,p^{\nu_b}\dd p$. Under $p\mapsto kp$, $v\mapsto k^5v$, $z\mapsto k^4z$ the form $p^m\re^{-\Wv(p)/z}\dd p$ is multiplied by $k^{m+1}$. The critical points of $\Wv$ are the $p$ with $p^5=v$, with the critical values $\frac54p^4$; colour $i$ corresponds to the critical point $p=\zeta^{-i}\varsigma$, $\varsigma^5=v$, where $\re^{-\Wv/z}=\re^{U_i/z}$ with $U_i\coloneqq-\frac54\zeta^i\varsigma^4$. Other normalisations of the model are compared in \cref{sec:sgn}.
\end{defn}

\begin{prop}[The model system and the T-system]
\label{prop:set-modrows}
\begin{itemize}
\item[(a)] If the symbols $I_m$ satisfy (Mod1) and (Mod2), the rows satisfy the \emph{model system}
\[
z\,\partial_vY_b=b^{\mathrm m}_b\,Y_{b+1}\quad(b\in\bbZ/5),\qquad b^{\mathrm m}=(b^{\mathrm m}_0,\dots,b^{\mathrm m}_4)\coloneqq\Bigl(1,-\frac1v,1,-1,-1\Bigr).
\]
\item[(b)] Fix a colour $i$, let $\varsigma$ be a formal fifth root of $v$, with $\partial_v\varsigma=1/(5\varsigma^4)$, and put $t\coloneqq z/\varsigma^4$; then $\partial_vU_i=-\zeta^i\varsigma^{-1}$ and $\partial_vt=-\frac45t/v$. A \emph{WKB column of colour $i$} is a column $\re^{U_i/z}y$ with $y\in\bbQ(\zeta)[\varsigma^{\pm1}]\fps z^5$, $\re^{U_i/z}$ being a formal exponential with $\partial_v\re^{U_i/z}=(\partial_vU_i/z)\re^{U_i/z}$; it is \emph{homogeneous} if $y_b=\varsigma^{\nu_b-1}F_b(t)$ with $F_b\in\bbQ(\zeta)\fps t$. A homogeneous WKB column solves the model system if and only if
\begin{gather*}
F_1=\mathcal T^{(i)}_{-1}F_0,\quad F_2=-\mathcal T^{(i)}_{-2}F_1,\quad F_3=\mathcal T^{(i)}_2F_2,\quad F_4=-\mathcal T^{(i)}_1F_3,\quad F_0=-\mathcal T^{(i)}_0F_4,\\
\mathcal T^{(i)}_k\coloneqq-\zeta^i+\frac t5\Bigl(k-4t\frac{\dd}{\dd t}\Bigr).
\end{gather*}
This T-system has a unique solution $F^{(i)}$ with $F^{(i)}_0(0)=1$, and $F^{(i)}_b(0)=\varepsilon_b\zeta^{ib}$. For colour $0$ we write $\mathcal T_k\coloneqq\mathcal T^{(0)}_k$ and $\Fs_b\coloneqq F^{(0)}_b$; thus $\Fs_b(0)=\varepsilon_b$, and for instance $\Fs_0=1-\frac3{20}t+\frac9{800}t^2+O(t^3)$. Two homogeneous WKB solutions of colour $i$ differ by a constant factor; and if $Y\neq0$ is a homogeneous WKB solution of colour $i$ and $d(z)\in\bbQ(\zeta)\fps z$, then $d(z)Y$ is a WKB solution of colour $i$, which is homogeneous only if $d(z)$ is constant.
\end{itemize}
\end{prop}

\begin{proof}
(a) By (Mod1), $z\partial_vY_0=Y_1$, $z\partial_vY_2=I_2=Y_3$, $z\partial_vY_3=-Y_4$, $z\partial_vY_4=-Y_0$, and with (Mod2) at $m=0$, $z\partial_vY_1=I_{-2}=v^{-1}I_3=-v^{-1}Y_2$.

(b) With $\partial_v\varsigma^c=\frac c5\varsigma^{c-5}$, for every $c$ and $F\in\bbQ(\zeta)\fps t$,
\[
z\,\partial_v\bigl[\re^{U_i/z}\varsigma^cF(t)\bigr]=\re^{U_i/z}\varsigma^{c-1}\Bigl(-\zeta^iF+\frac t5\bigl(cF-4tF'\bigr)\Bigr)=\re^{U_i/z}\varsigma^{c-1}\,\mathcal T^{(i)}_cF .
\]
So $z\partial_vY_b=\re^{U_i/z}\varsigma^{\nu_b-2}\mathcal T^{(i)}_{\nu_b-1}F_b$ and $b^{\mathrm m}_bY_{b+1}=b^{\mathrm m}_b\re^{U_i/z}\varsigma^{\nu_{b+1}-1}F_{b+1}$; for $b\neq1$, $\nu_{b+1}-1=\nu_b-2$, and $b^{\mathrm m}_1\varsigma^{\nu_2-1}=-\varsigma^{\nu_1-2}$, so the powers of $\varsigma$ match, and $(\nu_b-1)_b=(-1,-2,2,1,0)$ with the signs of $b^{\mathrm m}$ give the T-system.

On monomials, $\mathcal T^{(i)}_kt^n=-\zeta^it^n+\frac{k-4n}5t^{n+1}$. The first four relations define $F_1,\dots,F_4$ from $F_0$, and the closing relation is $F_0=-\mathcal T^{(i)}_0\mathcal T^{(i)}_1\mathcal T^{(i)}_2\mathcal T^{(i)}_{-2}\mathcal T^{(i)}_{-1}F_0$, whose operator maps $t^n$ to $t^n+4n\zeta^{4i}t^{n+1}+O(t^{n+2})$, as $(-\zeta^i)^5=-1$ and the five indices have sum $0$. So the coefficient of $t^{n+1}$ of the closing relation is $4n\zeta^{4i}\coef{t^n}F_0$ plus a combination of the $\coef{t^m}F_0$, $m<n$; the pivot is non-zero for $n\geq1$, so the solutions form a line, determined by $F_0(0)$, and two homogeneous WKB solutions of colour $i$ differ by a constant factor. At $t=0$, $\mathcal T^{(i)}_k$ acts as $-\zeta^i$, which gives $F^{(i)}_b(0)=\varepsilon_b\zeta^{ib}$. Since $z$ is a constant for $\partial_v$, $d(z)Y$ solves the model system. Write $d(z)=\sum_kd_k\varsigma^{4k}t^k$ and $Y_0=\re^{U_i/z}\varsigma^{-1}F_0$ with $F_0(0)\neq0$; if $d_1=\dots=d_{k-1}=0$, the coefficient of $t^k$ in $d(z)F_0$ is $d_0\coef{t^k}F_0+d_k\varsigma^{4k}F_0(0)$, free of $\varsigma$ only if $d_k=0$; by induction, $d(z)Y$ is homogeneous only if $d$ is constant.
\end{proof}

\input{sections/s02_setup_7}



